% ==========================================
% DOCUMENTO MAESTRO: APUNTES DE DERECHO
% PROCESAL CIVIL Y COMERCIAL
% Agosto 2026
% ==========================================

\documentclass[12pt,oneside]{book}

% ==========================================
% PAQUETES OBLIGATORIOS
% ==========================================

\usepackage[utf8]{inputenc}
\usepackage[T1]{fontenc}
\usepackage[spanish,es-nodecimaldot]{babel}

\usepackage[
    top=2.5cm,
    bottom=2.5cm,
    left=3cm,
    right=2.5cm,
    headheight=15pt
]{geometry}

\usepackage{amsmath}
\usepackage{amssymb}
\usepackage{mathtools}

\usepackage{graphicx}
\usepackage{float}
\usepackage{caption}
\usepackage{subcaption}

\usepackage{booktabs}
\usepackage{array}
\usepackage{longtable}
\usepackage{tabularx}

\usepackage{enumitem}

\usepackage{xcolor}
\usepackage[most]{tcolorbox}

\usepackage{fancyhdr}
\usepackage{ragged2e}

% ==========================================
% HIPERVÍNCULOS Y METADATOS PDF
% ==========================================

\usepackage[
    colorlinks=true,
    linkcolor=blue!50!black,
    citecolor=green!40!black,
    urlcolor=blue!60!black,
    pdftitle={Apuntes de Derecho Procesal Civil y Comercial},
    pdfauthor={Isaac Edgar Camacho Ocampo},
    pdfsubject={Apuntes universitarios},
    bookmarks=true
]{hyperref}

% ==========================================
% RUTAS DE IMÁGENES
% ==========================================

\graphicspath{
    {./img/}
    {./graficos/}
    {./figuras/}
}

% ==========================================
% CONFIGURACIÓN GENERAL DEL DOCUMENTO
% ==========================================

\setlength{\parindent}{0pt}
\setlength{\parskip}{0.5em}

\setlist[itemize]{
    topsep=0.3em,
    itemsep=0.2em
}

\setlist[enumerate]{
    topsep=0.3em,
    itemsep=0.2em
}

\clubpenalty=10000
\widowpenalty=10000

\setcounter{tocdepth}{2}

% ==========================================
% DEFINICIONES DE COMANDOS PERSONALIZADOS
% ==========================================

\newcommand{\abs}[1]{\left|#1\right|}
\newcommand{\norm}[1]{\left\lVert#1\right\rVert}
\newcommand{\vect}[1]{\mathbf{#1}}
\newcommand{\deriv}[2]{\frac{d#1}{d#2}}
\newcommand{\pderiv}[2]{\frac{\partial#1}{\partial#2}}

% Metadatos del documento
\newcommand{\universidad}{Facultad de Derecho}
\newcommand{\facultad}{Universidad de Buenos Aires}
\newcommand{\carrera}{Derecho}
\newcommand{\materia}{Derecho Procesal Civil y Comercial}
\newcommand{\titulo}{El Proceso de Conocimiento: Estructura, Demanda y Notificaciones}
\newcommand{\autor}{Isaac Edgar Camacho Ocampo}
\newcommand{\anio}{2026}

% ==========================================
% DEFINICIÓN DE CAJAS ACADÉMICAS
% ==========================================

\newtcolorbox{definition}[1]{
    enhanced,
    breakable,
    title={Definición: #1},
    fonttitle=\bfseries,
    colback=blue!3,
    colframe=blue!50!black,
    boxrule=0.8pt,
    arc=2mm
}

\newtcolorbox{important}{
    enhanced,
    breakable,
    title={Importante},
    fonttitle=\bfseries,
    colback=yellow!8,
    colframe=orange!70!black,
    boxrule=0.8pt,
    arc=2mm
}

\newtcolorbox{example}[1]{
    enhanced,
    breakable,
    title={Ejemplo: #1},
    fonttitle=\bfseries,
    colback=green!4,
    colframe=green!40!black,
    boxrule=0.8pt,
    arc=2mm
}

\newtcolorbox{formula}[1]{
    enhanced,
    breakable,
    title={Fórmula / Regla: #1},
    fonttitle=\bfseries,
    colback=purple!4,
    colframe=purple!50!black,
    boxrule=0.8pt,
    arc=2mm
}

\newtcolorbox{exercise}[1]{
    enhanced,
    breakable,
    title={Ejercicio: #1},
    fonttitle=\bfseries,
    colback=gray!5,
    colframe=gray!60!black,
    boxrule=0.8pt,
    arc=2mm
}

\newtcolorbox{note}{
    enhanced,
    breakable,
    title={Nota},
    fonttitle=\bfseries,
    colback=white,
    colframe=gray!50,
    boxrule=0.6pt,
    arc=2mm
}

% ==========================================
% ENCABEZADOS Y PIES DE PÁGINA
% ==========================================

\pagestyle{fancy}
\fancyhf{}

\fancyhead[L]{\nouppercase{\leftmark}}
\fancyhead[R]{\materia}

\fancyfoot[C]{\thepage}

\renewcommand{\headrulewidth}{0.4pt}
\renewcommand{\footrulewidth}{0pt}

\fancypagestyle{plain}{
    \fancyhf{}
    \fancyfoot[C]{\thepage}
    \renewcommand{\headrulewidth}{0pt}
}

% ==========================================
% INICIO DEL DOCUMENTO
% ==========================================

\begin{document}

% ==========================================
% PORTADA
% ==========================================

\begin{titlepage}

\thispagestyle{empty}

\begin{center}

\vspace*{1cm}

{\large \textsc{\facultad}}

\vspace{0.2cm}

{\large \textsc{\universidad}}

\vspace{0.2cm}

{\large \carrera}

\vspace{3cm}

{\Huge \textbf{\materia}}

\vspace{1cm}

{\LARGE \titulo}

\vspace{0.8cm}

\textit{La precisión en el lenguaje jurídico es el puente entre la norma abstracta y la justicia concreta.}

\vspace{1.2cm}

\rule{0.70\textwidth}{0.5pt}

\vspace{0.5cm}

{\large Apuntes de estudio}

\vspace{0.5cm}

\rule{0.70\textwidth}{0.5pt}

\vfill

{\large \autor}

\vspace{0.3cm}

{\large \anio}

\end{center}

\vfill

\begin{flushleft}
\textit{Este es un modesto aporte a los estudiantes de ``Derecho Procesal Civil y Comercial'' no pretende sustituir el material oficial de la materia.}
\end{flushleft}

\end{titlepage}

% ==========================================
% TABLA DE CONTENIDOS
% ==========================================

\tableofcontents

% ==========================================
% PRESENTACIÓN
% ==========================================

\chapter*{Presentación}
\addcontentsline{toc}{chapter}{Presentación}

Este material constituye una compilación de apuntes de clase correspondientes a Derecho Procesal Civil y Comercial, dictada en la Facultad de Derecho de la Universidad de Buenos Aires durante el mes de agosto de 2026.

La clase recorre el camino que va desde la norma jurídica abstracta hasta el expediente judicial concreto. Partiendo de una premisa fundamental ---el derecho se trabaja siempre con la palabra---, se construye una visión integrada de cómo operan los procesos judiciales, cómo se formulan las demandas y cuáles son las reglas que rigen la comunicación entre las partes y el tribunal.

Estos apuntes no pretenden reemplazar la bibliografía oficial de la materia, sino servir como herramienta complementaria de estudio y como referencia para organizar los conceptos desarrollados en clase.

% ==========================================
% TABLA RESUMEN DE TEMAS
% ==========================================

\chapter*{Temas Tratados en esta Clase}
\addcontentsline{toc}{chapter}{Temas Tratados}

A lo largo de esta clase se abordaron los siguientes ejes temáticos:

\begin{itemize}
    \item El lenguaje jurídico y la estructura lógica de las normas (mandato, permiso y sanción).
    \item Proceso de conocimiento vs. proceso ejecutivo: efectos y cosa juzgada.
    \item La construcción de la demanda: del cliente a la estrategia procesal.
    \item Inicio del proceso, competencia e incidentes.
    \item Taller práctico de redacción de escritos judiciales.
    \item Régimen de notificaciones, cómputo de plazos y preclusión procesal.
\end{itemize}

% ==========================================
% SECCIÓN: RELEVANCIA PROFESIONAL Y PEDAGÓGICA
% ==========================================

\chapter*{¿Para Qué Sirve Esto en la Vida Profesional?}
\addcontentsline{toc}{chapter}{Relevancia Profesional}

Todo abogado trabaja, en definitiva, con un único material: la palabra. Distinguir un proceso de conocimiento de uno ejecutivo, construir una demanda a partir del relato incompleto de un cliente y dominar los plazos y las notificaciones es la base técnica de cualquier estrategia procesal. Un error en el cómputo de un plazo o en la redacción de un escrito puede costarle un derecho al cliente: este contenido es la caja de herramientas con la que se litiga todos los días.

\section*{¿Es Útil para la Vida? ¿Por Qué?}

Más allá de lo profesional, esta clase enseña una forma de pensar. Distinguir lo que ocurrió de lo que se puede probar, comunicar con claridad y precisión, y entender que toda norma de convivencia implica una consecuencia son habilidades que exceden el tribunal y sirven para negociar, argumentar y decidir en cualquier ámbito de la vida.

% ==========================================
% ÍNDICE TEMÁTICO CONCEPTUAL
% ==========================================

\chapter*{Índice Temático Conceptual}
\addcontentsline{toc}{chapter}{Índice Temático}

\section*{Cómo se Relacionan los Bloques}

Pensemos en un semáforo mal cruzado. Primero existe una norma de convivencia (Bloque 1: ¿está prohibido u obligado cruzar?). Cuando el hecho ocurre y hay un choque, hace falta un procedimiento para resolverlo, y no todos los conflictos se tramitan igual (Bloque 2: ¿conocimiento o ejecutivo?). Para llegar a un juez, alguien tiene que contar lo que pasó y convertirlo en una estrategia (Bloque 3: la demanda). Ese relato se presenta formalmente y dispara efectos concretos, como abrir incidentes accesorios (Bloque 4). Todo eso se traduce en escritos con reglas propias (Bloque 5) y cada paso del expediente le llega a las partes según reglas de tiempo que, si se ignoran, hacen perder el derecho a reclamar (Bloque 6). Es la misma historia contada en seis capítulos: de la norma general al expediente concreto.

\section*{Bloques de Contenido}

\textbf{⚖ El Derecho Como Lenguaje (Unidades 1 a 3)}
\begin{itemize}
    \item Unidad 1: La palabra como herramienta de trabajo --- lenguaje jurídico, interpretación, aplicación.
    \item Unidad 2: Estructura lógica de las normas --- mandato, prohibición, permiso.
    \item Unidad 3: La sanción como elemento esencial --- sanción jurídica, social, disciplinaria.
\end{itemize}

\textbf{⚖ Tipos de Proceso Judicial (Unidades 4 a 8)}
\begin{itemize}
    \item Unidad 4: Organización del Código Procesal --- parte general, parte especial.
    \item Unidad 5: Proceso de conocimiento y proceso ejecutivo --- conocer los hechos vs. validez del título.
    \item Unidad 6: Cosa juzgada formal y material --- firmeza, revisión posterior.
    \item Unidad 7: Ordinario, sumario y sumarísimo --- celeridad, plazos, prueba.
    \item Unidad 8: Proceso de ejecución de sentencia --- cumplimiento forzado, etapa común.
\end{itemize}

\textbf{✎ La Construcción de la Demanda (Unidades 9 a 12)}
\begin{itemize}
    \item Unidad 9: De los hechos al proceso --- relato, escrito, prueba, derecho.
    \item Unidad 10: La entrevista con el cliente --- relato incompleto, interpretación profesional.
    \item Unidad 11: Estrategia procesal y selección de hechos --- relevancia jurídica, prueba disponible.
    \item Unidad 12: Verdad real y verdad procesal --- verdad formal, sentencia, reconstrucción.
\end{itemize}

\textbf{⚖ Inicio del Proceso, Competencia e Incidentes (Unidades 13 a 15)}
\begin{itemize}
    \item Unidad 13: Mediación prejudicial e inicio del proceso --- instancia previa, demanda.
    \item Unidad 14: Efectos de la interposición de la demanda --- hechos controvertidos, competencia.
    \item Unidad 15: Los incidentes --- cuestión accesoria, cuerda separada.
\end{itemize}

\textbf{✒ Taller de Redacción de Escritos (Unidades 16 a 18)}
\begin{itemize}
    \item Unidad 16: Estructura formal de un escrito judicial --- título, sumario, encabezamiento, petitorio.
    \item Unidad 17: Escrito de mero trámite --- apertura a prueba, sumario.
    \item Unidad 18: Escrito fundado --- justificación, fuerza mayor, prueba testimonial.
\end{itemize}

\textbf{⏱ Notificaciones y Plazos Procesales (Unidades 19 y 20)}
\begin{itemize}
    \item Unidad 19: Formas de notificación --- personal, por cédula, por nota.
    \item Unidad 20: Cómputo de plazos y preclusión --- día de nota, vencimiento, cierre del acto.
\end{itemize}

% ==========================================
% BLOQUE 1: EL DERECHO COMO LENGUAJE
% ==========================================

\chapter{El Derecho Como Lenguaje}

\section{Unidad 1: La Palabra Como Herramienta de Trabajo del Abogado}

Una pregunta fundamental abre esta unidad: ¿con qué material trabaja realmente un abogado? La respuesta es sencilla pero determinante: con la palabra.

Cuando estudiamos algo, creemos conocerlo, pero necesitamos poder verlo con claridad. En el caso del trabajo jurídico, esa claridad se alcanza solo mediante el lenguaje, sean enunciados escritos u orales. Las normas, en definitiva, constituyen proposiciones lingüísticas que deben ser interpretadas, aplicadas e invocadas. El material del derecho es siempre textual.

\begin{example}{La materia prima del abogado}
Consideremos un enunciado penal clásico: ``el que matare a otro tendrá pena de 15 a 25 años''. Frente a ese enunciado, la función del abogado es, en primer lugar, interpretarlo y, luego, aplicarlo o invocarlo. Pero siempre, en todos los casos, sobre la base de la palabra.
\end{example}

El trabajo del abogado no se relaciona con ``el universo'' ni con cuestiones ajenas a lo textual. La palabra ---hablada o escrita--- constituye el único instrumento verdadero de la profesión. Por ese motivo resulta crucial que los estudiantes se entrenen en expresarse de manera clara y precisa, tanto oralmente como por escrito, un ejercicio poco desarrollado en la formación universitaria tradicional.

\section{Unidad 2: Estructura Lógica de las Normas: Mandato y Permiso}

A partir de la idea de la norma como enunciado, puede pensarse en las normas jurídicas del mismo modo en que las primeras comunidades humanas habrían pensado sus normas de convivencia: como reglas que se ponen en ejercicio frente a un hecho concreto.

Juzgar una cuestión jurídica es confrontar un enunciado normativo con la realidad. Según se estudia en Introducción al Derecho y en Filosofía del Derecho, toda norma jurídica contiene en el fondo solo dos posibilidades lógicas: obligar o permitir.

\begin{example}{La prohibición traducida al mandato}
Consideremos una norma de tránsito: ``está prohibido estacionar''. Traduciendo esta norma de manera positiva, el sentido es exactamente: ``es obligatorio no estacionar''. La prohibición equivale, en sentido positivo, a una obligación. Lo mismo ocurre con el ejemplo penal ya mencionado: la norma que castiga al ``que matare'' con una pena de 15 a 25 años equivale a decir, en sentido positivo, que es obligatorio no matar ---bajo pena de sufrir esa consecuencia si se lo hace---.
\end{example}

\begin{important}
Toda norma jurídica cumple una de dos funciones: obligar (lo que incluye, en su formulación negativa, a las prohibiciones) o permitir. Cuando una norma no prohíbe ni obliga una conducta, sino que la autoriza, dicha norma es facultativa: deja al destinatario la posibilidad de actuar o no actuar, sin consecuencia alguna en ninguno de los dos casos.
\end{important}

\section{Unidad 3: La Sanción Como Elemento Esencial de Toda Norma}

Para que una norma sea verdaderamente una norma, debe llevar aparejada una sanción. Existen distintos tipos de sanciones: económicas (como la indemnización de daños), penales, disciplinarias y sociales, siendo estas últimas sanciones de hecho que no se traducen en una consecuencia formal, pero que cumplen la función estructural de toda sanción.

\begin{example}{La sanción social}
Imaginemos que un profesor llega a clase sin saludar a nadie, maltratando verbalmente a los presentes, o vistiendo un traje de color rojo con corbata verde y zapatos anaranjados. Probablemente la universidad no aplicaría ninguna sanción formal disciplinaria ni jurídica, pero sí existiría una reacción social clara: todos pensarían ``mal gusto'' o formarían un juicio negativo. Esa reacción silenciosa es, en sí misma, una sanción: una sanción social que cumple la misma función estructural que cualquier otra sanción.
\end{example}

\begin{important}
Toda norma implica una sanción. Cuando una norma no implica una sanción, en realidad ya no cumple ninguna función como tal. A partir de esta idea, se forma en los estudiantes una manera de pensar propiamente jurídica.
\end{important}

% ==========================================
% BLOQUE 2: TIPOS DE PROCESO JUDICIAL
% ==========================================

\chapter{Tipos de Proceso Judicial}

\section{Unidad 4: Organización del Código Procesal}

El Código Procesal se organiza en dos grandes divisiones: una \textbf{parte general} y una \textbf{parte especial}. La parte general regula lo relativo al órgano judicial y resulta aplicable a todos los tipos de proceso. La parte especial, en cambio, contiene la regulación específica de cada tipo procesal.

Existen dos grandes tipos de proceso: el \textbf{ordinario} y el \textbf{ejecutivo}, distinción que es fundamental para toda la teoría procesal.

\section{Unidad 5: Proceso de Conocimiento y Proceso Ejecutivo}

El \textbf{proceso de conocimiento} se caracteriza por la amplia capacidad probatoria de las partes, las amplias facultades del juez para conocer los hechos y ordenar prueba, y las amplias facultades de las partes para argumentar. Su objeto es conocer los hechos, es decir, la versión de lo ocurrido que cada parte aporta al proceso.

\begin{important}
En el proceso de conocimiento no se conocen los hechos tal como ocurrieron en la realidad. Lo que se conoce es la versión de los hechos que cada parte relata y cree que sucedió. El objetivo es reconstruir de la mejor manera posible esa versión, aplicar el derecho correspondiente y arribar a un resultado: la sentencia.
\end{important}

Frente al proceso de conocimiento, el \textbf{proceso ejecutivo} opera sobre presupuesto totalmente distinto: la existencia de un título ejecutivo. En este tipo de proceso, lo que el juez examina no son los hechos que dieron origen al título, sino su \textbf{validez formal}: la existencia de una fecha, una firma y los demás requisitos que la ley exige para cada tipo de instrumento.

\begin{example}{El cheque sin fondos}
Un vendedor le vende a un comprador una sustancia y recibe como pago un cheque que luego resulta sin fondos. El vendedor puede iniciar la ejecución del cheque. El título no expresa la causa de la obligación; el juez únicamente examinará su validez formal, sin indagar el negocio subyacente. ``El título no va a decir'' cuál fue la causa de la obligación; esa es otra cuestión, ajena al proceso ejecutivo.
\end{example}

Los títulos ejecutivos existen para agilizar la circulación comercial y otorgar seguridad jurídica, evitando que cada operación deba resolverse mediante un proceso de conocimiento ---más lento y trabajoso--- en caso de incumplimiento.

\begin{example}{Los quintales de trigo con gorgojo}
Una persona compra a un vendedor cinco quintales de trigo y paga con un cheque. Al recibir la mercadería, descubre que está infestada de gorgojo y solo puede recuperarse un quintal. Esa persona no puede negarse a pagar el cheque, porque el título es válido y el vendedor puede ejecutarlo igualmente. Lo que sí podrá hacer, con posterioridad, es iniciar un proceso de conocimiento contra el vendedor para reclamar por el incumplimiento del contrato de compraventa, ya que se trata de una causa distinta a la validez del título.
\end{example}

\section{Unidad 6: Cosa Juzgada Formal y Cosa Juzgada Material}

En el proceso de conocimiento, la sentencia produce el efecto de \textbf{cosa juzgada}: una vez firme, la cuestión no puede volver a discutirse. En el proceso ejecutivo, en cambio, aunque también existe una sentencia que ordena el pago, la cosa juzgada que allí se produce es meramente formal.

\begin{important}
La sentencia dictada en el proceso ejecutivo hace cosa juzgada formal respecto del título, pero no impide que, posteriormente, la parte condenada inicie un proceso de conocimiento para discutir la causa del negocio subyacente ---por ejemplo, para reclamar por la mercadería defectuosa recibida---. En cambio, la sentencia dictada en un proceso de conocimiento produce cosa juzgada material: agota definitivamente la controversia y no puede volver a discutirse en ningún proceso posterior.
\end{important}

\section{Unidad 7: Ordinario, Sumario y Sumarísimo}

Existen varios tipos de proceso de conocimiento, entre ellos el \textbf{sumario} y el \textbf{sumarísimo}, además del \textbf{ordinario}. Si bien en la Capital Federal el proceso sumario se encuentra derogado en la práctica, el Código todavía conserva referencias a él, y en varias provincias argentinas ---especialmente en la provincia de Buenos Aires--- continúa vigente.

Todos estos procesos son procesos de conocimiento. La diferencia entre ellos radica en la \textbf{celeridad}: el sumario y el sumarísimo tienen plazos más cortos, menor cantidad de prueba admisible y menos recursos disponibles, priorizando la velocidad por sobre la extensión del debate. En todos los casos, lo que se busca conocer son los hechos, a diferencia del proceso ejecutivo, en el que lo que se busca conocer es la validez del título.

\section{Unidad 8: El Proceso de Ejecución de Sentencia}

Tanto el proceso de conocimiento como el proceso ejecutivo concluyen con una sentencia de condena. En ambos casos, si la parte condenada no cumple voluntariamente dentro del plazo fijado, corresponde iniciar el \textbf{proceso de ejecución de sentencia}.

\begin{example}{Los 100 pesos y los 30 días}
Una sentencia hipotética hace lugar a una demanda y condena al demandado a pagar 100 pesos en un plazo de 30 días. Si transcurrido ese plazo el demandado no paga, corresponde iniciar el proceso de ejecución de sentencia, exactamente del mismo modo que ocurriría si, en un juicio ejecutivo, se ordenara pagar un cheque y transcurrieran los 10 días previstos sin que el pago se efectuara.
\end{example}

El proceso de ejecución va a ser común a ambos tipos de proceso, porque lo que se busca es hacer efectivo lo que el juez mandó. La diferencia entre ambos procesos reside únicamente en el camino previo: en un caso, la decisión se alcanzó tras el estudio de los hechos y la producción de pruebas; en el otro, tras el examen de un documento que reunía los requisitos formales exigidos por la ley. Pero, una vez dictada la sentencia, si las partes no cumplen voluntariamente, la etapa de ejecución es idéntica en su lógica para ambos tipos de proceso.

% ==========================================
% BLOQUE 3: CONSTRUCCIÓN DE LA DEMANDA
% ==========================================

\chapter{La Construcción de la Demanda}

\section{Unidad 9: De los Hechos al Proceso}

Un conflicto generado por hechos concretos ocurridos entre dos personas llega a convertirse en un proceso judicial a través de la \textbf{demanda}, entendida como un escrito acompañado de prueba.

La demanda integra varios elementos: en primer lugar, el relato de los hechos aportado por el cliente; en segundo lugar, la prueba que respalda ese relato; y, finalmente, el aporte propio del abogado, que consiste en la \textbf{fundamentación jurídica} ---el encuadre de esos hechos dentro del derecho aplicable---. El objetivo de todo este proceso de construcción es que el juez tenga por establecidos los hechos tal como los presenta el cliente, y ese es, en definitiva, el camino para obtener una sentencia favorable: convencer al juez.

\section{Unidad 10: La Entrevista con el Cliente y la Reconstrucción del Relato}

Para entender las limitaciones de cualquier relato, consideremos un ejemplo: si se pide a alguien que cuente en pocos segundos cómo llegó a un lugar, esa persona relata una síntesis que, por mejor que sea su esfuerzo, resulta necesariamente incompleta.

\begin{example}{Los 200 millones de sensaciones}
Imaginemos una persona que cuenta cómo llegó a la facultad. Relata que se preparó, tomó un colectivo durante unos 40 minutos y luego hizo una combinación de subte entre dos líneas. Sin embargo, durante esos 40 minutos de trayecto, esa persona percibió ---consciente o inconscientemente--- una enorme cantidad de estímulos: personas que pasaron, un vendedor ambulante, alguien que la empujó, entre muchísimos otros detalles que ni siquiera llegó a registrar. Por mejor que fuera el esfuerzo de síntesis, el relato resultaba necesariamente incompleto.
\end{example}

Lo mismo ocurre con el relato de cualquier cliente: por más detallado que sea, jamás reproduce la totalidad de lo ocurrido. El trabajo del abogado frente al relato del cliente no es transcribirlo literalmente, sino \textbf{reconstruirlo con criterio profesional}: elaborar una estrategia y organizar la narración en función de ella.

El cliente, incluso con la mejor voluntad, relata los hechos desde su propio punto de vista y con su propia carga emocional, por lo que rara vez cuenta todo lo ocurrido o lo cuenta de manera completamente objetiva. Por ello, la entrevista con el cliente es instancia central del trabajo profesional: en ella deben formularse preguntas bien diseñadas y desplegarse la capacidad de interpretar relatos de personas de muy distintos niveles socioculturales, para obtener el máximo de información posible.

\section{Unidad 11: Estrategia Procesal y Selección de los Hechos}

¿Qué elementos determinan que un abogado decida incluir o no un hecho determinado en su relato estratégico? Tres criterios centrales emergen de esta pregunta:

\begin{enumerate}
    \item La \textbf{relevancia jurídica} del hecho: si favorece o no la posición del cliente.
    \item La \textbf{calificación jurídica} adecuada del hecho.
    \item La \textbf{posibilidad de probarlo}: sobre todo, este tercer criterio resulta decisivo.
\end{enumerate}

\begin{important}
Todo aquello que se afirma y no se puede probar juega en contra. Si un hecho es relevante o importante pero no existen elementos para acreditarlo, lo estratégicamente correcto es omitirlo o dejarlo en un segundo plano. El abogado debe pensar de manera estratégica, anticipando además cuál será la posición de la contraparte, para no dejar espacios ni flancos débiles en el propio relato.
\end{important}

Es importante recordar que cuando una demanda llega al juzgado, habitualmente no la lee primero el juez, sino personal del juzgado con múltiples preocupaciones y expedientes pendientes. Por ese motivo, los escritos judicales deben ser concretos, precisos y ordenados, cualidades que no son sencillas de lograr y requieren un trabajo genuino de síntesis.

\section{Unidad 12: Verdad Real y Verdad Procesal}

Una pregunta fundamental surje al estudiar la sentencia: ¿coincide exactamente lo que establece una sentencia con lo que ocurrió en la realidad? En el mejor de los casos, ambas cosas se parecen, pero no son idénticas.

\begin{example}{La verdad que nunca se conoce del todo}
Existe una verdad de lo que realmente ocurrió, a la que ni siquiera los propios abogados o el juez tienen acceso directo. Lo que sí existe es un conjunto de pruebas y una contraposición de relatos entre las partes, de cuya confrontación surge un resultado: la sentencia, que constituye la llamada verdad formal o verdad jurídica. Esa verdad formal es la que finalmente se impone sobre hechos que, muchas veces, ocurrieron años atrás y que jamás podrán reconstruirse con exactitud absoluta.
\end{example}

% ==========================================
% BLOQUE 4: INICIO PROCESO, COMPETENCIA, INCIDENTES
% ==========================================

\chapter{Inicio del Proceso, Competencia e Incidentes}

\section{Unidad 13: Mediación Prejudicial e Inicio del Proceso}

¿Cuál es el punto de inicio del proceso judicial? La mediación prejudicial, cuando es exigida, constituye un requisito previo para ciertos tipos de juicio ---la mayoría de las provincias argentinas la exigen actualmente con carácter obligatorio---. Sin embargo, la mediación es \textbf{previa} al proceso: no forma parte del proceso propiamente dicho.

\begin{important}
El proceso judicial en sentido técnico se inicia recién con la interposición de la demanda. La mediación prejudicial, si bien es una etapa ineludible en muchas jurisdicciones, constituye una etapa anterior al proceso, no su comienzo.
\end{important}

\section{Unidad 14: Efectos de la Interposición de la Demanda}

La interposición de la demanda produce dos efectos fundamentales que, con seguridad, serán tema de examen: la fijación de los \textbf{hechos controvertidos} y la fijación de la \textbf{competencia}.

\begin{example}{El accidente entre San Isidro y Santa Fe}
Un vehículo, cuyo titular vive en la Ciudad Autónoma de Buenos Aires, choca contra otro vehículo cuya compañía aseguradora tiene su sede en la provincia de Santa Fe. El hecho ocurrió en la localidad de San Isidro. El ordenamiento concede al actor la posibilidad de elegir entre distintas opciones de competencia ---el domicilio del demandado, el lugar del hecho, o el domicilio de la aseguradora, según el caso---. Esa elección queda definitivamente fijada una vez presentada la demanda, sin que el actor pueda disponer libremente de ella con posterioridad.
\end{example}

Si bien el actor no dispone en general de la competencia, el propio ordenamiento procesal, en ciertos supuestos, habilita al actor a elegir entre distintas opciones, y en esa elección el abogado define cuál es la alternativa más conveniente para su cliente.

\section{Unidad 15: Los Incidentes}

Un \textbf{incidente} es un proceso secundario, vinculado al proceso principal, que resuelve una cuestión accesoria ---no definitoria del fondo del asunto--- pero que incide en el avance del proceso principal.

\begin{example}{La excepción de incompetencia}
Un actor inicia una demanda por incumplimiento de contrato ante determinado juez, y el demandado, al contestar la demanda, plantea que ese juez es incompetente y que correspondería la intervención de otro. Entre el actor y el demandado se genera una controversia que deberá resolverse antes de continuar. Si el demandado gana ese planteo, no gana el juicio: gana únicamente esa controversia parcial, denominada incidente, que incide en el proceso pero no lo define en cuanto al fondo.
\end{example}

\begin{example}{El beneficio de litigar sin gastos}
Cuando se inicia un juicio de contenido patrimonial y el actor no cuenta con los fondos suficientes para afrontar las costas del proceso, puede presentar un incidente de beneficio de litigar sin gastos. Su objetivo es permitir el acceso a la justicia a quien no dispone de recursos económicos, eximiéndolo, en principio, del pago de las tasas de justicia iniciales y, eventualmente, de las costas al finalizar el proceso, si el incidente resulta concedido.
\end{example}

En épocas en que los expedientes se tramitaban en papel, el incidente se sustanciaba en un expediente separado ---el llamado ``expediente por cuerda separada''--- físicamente atado con un hilo al expediente principal, de donde proviene la expresión ``ir por cuerda''. Actualmente, el juez puede optar por no formar un expediente separado y tramitar el incidente dentro del expediente principal; esa discrecionalidad sigue siendo válida hoy.

% ==========================================
% BLOQUE 5: TALLER REDACCIÓN ESCRITOS
% ==========================================

\chapter{Taller de Redacción de Escritos Judiciales}

\section{Unidad 16: Estructura Formal de un Escrito Judicial}

Todo escrito judicial debe contener, en orden, los siguientes elementos:

\begin{important}
1. \textbf{Título o sumario}: un resumen ínfimo, en pocas palabras, de lo que se solicita (por ejemplo: ``Solicita apertura a prueba'').

2. \textbf{Encabezamiento}: a quién se dirige el escrito (``Señor Juez'') y quién se presenta ---el abogado con el patrocinio o apoderamiento correspondiente, o la parte por su propio derecho---.

3. \textbf{Individualización de la causa}: los autos caratulados y el número de expediente.

4. \textbf{Exposición y petitorio}: el desarrollo de lo que se solicita y el pedido concreto dirigido al juez.
\end{important}

El domicilio real y el Documento Nacional de Identidad (DNI) del presentante se consignan únicamente en el \textbf{primer escrito} del expediente ---usualmente la demanda o la contestación---, momento en el que también se constituye el domicilio procesal para todas las notificaciones sucesivas. En los escritos posteriores esa información ya no vuelve a repetirse.

\section{Unidad 17: Ejercicio de Redacción: Escrito de Mero Trámite}

Un escrito de \textbf{mero trámite} es un escrito simple, en el que la parte se limita a informar el estado de los autos y a formular un pedido puntual, sin necesidad de una argumentación extensa.

\begin{example}{Escrito modelo redactado en clase}
\textit{``Solicita apertura a prueba. Señor Juez: Fulano de Tal, por derecho propio, con domicilio constituido en los autos caratulados `[carátula]', expediente n.° [número], a Vuestra Señoría respetuosamente digo: Que vengo por el presente, conforme al estado de autos, a solicitar se abra la causa a prueba. Proveer de conformidad. Será justicia.''}
\end{example}

\section{Unidad 18: Ejercicio de Redacción: Escrito Fundado}

Un escrito \textbf{fundado} requiere mayor desarrollo argumentativo. Consideremos una situación: un testigo clave de la parte debe viajar en la fecha fijada para la audiencia testimonial, por razones laborales o médicas, y no podrá comparecer. Deben redactarse un escrito que justifique esa ausencia y solicite el cambio de fecha.

\begin{example}{Primer intento: un escrito insuficiente}
Un primer estudiante propone un escrito breve que se limita a manifestar que el testigo no podría presentarse en la audiencia fijada por cuestiones laborales, hasta finalizar el mes de enero, y solicita que la nueva audiencia no se fije antes de esa fecha.
\end{example}

El escrito es válido en su estructura, pero puede perfeccionarse. Para mejorarlo, pueden incorporarse dos elementos adicionales:

\textbf{Primero}, calificar expresamente la situación del testigo como un caso de \textbf{fuerza mayor} ---por ejemplo, la imposibilidad de cambiar pasajes ya adquiridos---.

\textbf{Segundo}, y sobre todo, explicar por qué ese testigo resulta determinante o trascendente para la causa. Si el escrito no explica la relevancia del testigo, se corre el riesgo de que el juez deniegue el pedido de cambio de audiencia por considerarlo poco fundado.

\begin{important}
No alcanza con solicitar algo al juez; es necesario fundamentar por qué se lo solicita. Cuanto mayor es la carga argumentativa que exige el pedido ---como en este caso, donde se solicita reprogramar una audiencia ya fijada---, mayor debe ser el esfuerzo de justificación del escrito, explicando tanto la causa de fuerza mayor como la relevancia concreta de la prueba que se busca preservar.
\end{important}

% ==========================================
% BLOQUE 6: NOTIFICACIONES Y PLAZOS
% ==========================================

\chapter{Régimen de Notificaciones y Plazos Procesales}

\section{Unidad 19: Formas de Notificación}

Notificar es, en esencia, comunicar una resolución judicial. Toda resolución ---sea un traslado de demanda, una providencia simple o una sentencia--- debe ser notificada de alguna de las formas previstas por el Código.

Se identifican tres grandes formas de notificación:

\begin{enumerate}
    \item \textbf{Notificación personal}: se produce, típicamente, cuando la parte o su letrado concurren a tomar vista del expediente y quedan notificados de su contenido en ese acto. Una estrategia habitual de los actores es concurrir personalmente a notificarse cuando el demandado se muestra poco interesado en hacerlo voluntariamente, precisamente para no quedar a la espera de la diligencia del otro y poder continuar impulsando el proceso.
    
    \item \textbf{Notificación por cédula}: se trata de un formulario específico que puede confeccionarse en soporte papel o digital, siendo hoy mayoritariamente digital. Existen excepciones que aún requieren soporte papel, como el traslado de la demanda cuando el demandado todavía no constituyó domicilio procesal en el expediente.
    
    \item \textbf{Notificación por nota o ministerio de la ley}: el principio general es que toda resolución judicial queda notificada los días martes y viernes automáticamente, excepto que se trate de alguna de las resoluciones que, conforme al Artículo 133 del Código Procesal Civil y Comercial de la Nación, deben notificarse necesariamente por cédula o de manera personal. Si el martes o el viernes correspondiente fuera feriado, la notificación se traslada al primer día hábil siguiente.
\end{enumerate}

Las dos primeras formas ---personal y por cédula--- cumplen, en la práctica, una función equivalente.

\section{Unidad 20: Cómputo de Plazos y Preclusión Procesal}

El cómputo de los plazos procesales requiere especial cuidado. Si una resolución es dictada el martes, por ejemplo, no queda notificada ese mismo día: el día de salida de la resolución nunca corre a los fines del cómputo. Como el martes ya pasó, la resolución quedaría notificada recién el viernes siguiente, que es el próximo día de nota disponible.

\begin{example}{El cómputo del plazo: nunca corre el día en que sale la resolución}
Una resolución dictada el martes 18 de agosto no queda notificada ese mismo día. Como el siguiente día de nota es el viernes 21, la notificación se produce recién ese viernes, y el plazo para responder comienza a correr a partir del primer día hábil siguiente a la notificación ---es decir, el lunes posterior---, nunca desde el día en que la resolución fue dictada ni desde el día mismo de la notificación.
\end{example}

Un error en el cómputo del plazo puede hacer que opere la \textbf{preclusión}, concepto que designa el cierre definitivo de cada acto procesal una vez vencido su plazo, para permitir que el proceso avance a la etapa siguiente.

\begin{important}
Todo acto procesal debe cerrarse para que el proceso pueda continuar. La preclusión opera de dos maneras: si la parte cumple con el acto dentro del plazo otorgado (por ejemplo, contesta la demanda dentro de los 15 días previstos), el acto se cierra por cumplimiento, y no puede completarse ni ampliarse después, aun si el plazo total todavía no venció. Si, en cambio, la parte deja transcurrir el plazo sin cumplir, el acto se cierra por vencimiento. En ambos casos, una vez cerrado el acto, no puede volver a ejercerse.

Como excepción, la demanda es el único escrito que puede ampliarse con posterioridad a su presentación, pero únicamente \textbf{antes de que se haya notificado al demandado}, ya que hasta ese momento la relación procesal todavía no involucra a la contraparte.
\end{important}

% ==========================================
% CUADRO CORNELL
% ==========================================

\chapter{Cuadro de Repaso: Método Cornell}

A continuación se presenta un cuadro de estudio siguiendo el método Cornell, que organiza la información en preguntas clave, sus respuestas correspondientes y un resumen final de síntesis.

\begin{table}[H]
\centering
\caption*{\textbf{Cuadro Cornell --- Apuntes de Derecho Procesal Civil y Comercial}}
\vspace{0.3cm}
\begin{tabularx}{\textwidth}{
    >{\raggedright\arraybackslash}p{0.28\textwidth}
    >{\raggedright\arraybackslash}X
}
\toprule
\textbf{Preguntas / Conceptos Clave} & \textbf{Notas / Respuestas} \\
\midrule

\textbf{¿Con qué material trabaja siempre el abogado?} & 
Con la palabra, escrita u oral. Las normas son enunciados que deben interpretarse y aplicarse; no hay otro instrumento de trabajo posible en la práctica jurídica. \\

\textbf{¿Qué dos funciones puede cumplir toda norma jurídica?} & 
Obligar (incluida su forma negativa, la prohibición) o permitir. Toda norma redactada de cualquier modo se reduce siempre a uno de estos dos mandatos. \\

\textbf{¿Por qué toda norma necesita una sanción?} & 
Porque sin sanción una norma deja de cumplir su función como tal. Las sanciones pueden ser jurídicas, económicas, disciplinarias o sociales. \\

\textbf{¿Qué diferencia al proceso de conocimiento del proceso ejecutivo?} & 
El de conocimiento busca establecer los hechos mediante amplia prueba y debate; el ejecutivo solo examina la validez formal de un título (fecha, firma, requisitos), sin indagar la causa del negocio. \\

\textbf{¿Qué diferencia hay entre cosa juzgada formal y material?} & 
La cosa juzgada formal (propia del proceso ejecutivo) permite discutir después la causa en un proceso de conocimiento; la cosa juzgada material (propia del proceso de conocimiento) agota definitivamente la controversia. \\

\textbf{¿En qué se diferencian el ordinario, el sumario y el sumarísimo?} & 
Los tres son procesos de conocimiento; se diferencian en la celeridad: plazos, cantidad de prueba admisible y recursos disponibles. El sumario está derogado en la Capital Federal pero vigente en varias provincias. \\

\textbf{¿Cuándo corresponde el proceso de ejecución de sentencia?} & 
Cuando, vencido el plazo fijado en la sentencia (de conocimiento o ejecutiva), la parte condenada no cumple voluntariamente. Es una etapa común a ambos tipos de proceso. \\

\textbf{¿Por qué el relato de un cliente siempre es incompleto?} & 
Porque nadie puede reconstruir la totalidad de lo percibido en una situación. El abogado debe reconstruir ese relato con criterio profesional y estratégico. \\

\textbf{¿Qué criterios definen qué hechos incluir en una demanda?} & 
La relevancia jurídica del hecho, su correcta calificación legal y, sobre todo, la posibilidad concreta de probarlo. Lo que no se puede probar, mejor omitirlo. \\

\textbf{¿La sentencia refleja la verdad de lo que realmente ocurrió?} & 
No exactamente. La sentencia establece una verdad formal o jurídica, construida a partir de relatos, pruebas y contraposiciones entre partes; la verdad real de los hechos nunca se conoce con certeza absoluta. \\

\textbf{¿Con qué acto se inicia formalmente el proceso judicial?} & 
Con la demanda. La mediación prejudicial, cuando es exigida, es un requisito previo, pero no forma parte del proceso judicial propiamente dicho. \\

\textbf{¿Qué efectos produce la interposición de la demanda?} & 
Fija los hechos controvertidos sobre los que versará el proceso y fija la competencia del juez interviniente, dentro de las opciones que la ley concede al actor. \\

\textbf{¿Qué es un incidente y qué ejemplos existen?} & 
Un proceso accesorio vinculado al principal que resuelve una cuestión que incide en su avance sin definir el fondo. Ejemplos: la excepción de incompetencia y el beneficio de litigar sin gastos. \\

\textbf{¿Qué partes no pueden faltar en un escrito judicial?} & 
Título o sumario, encabezamiento (a quién se dirige y quién se presenta), individualización de la causa (carátula y número de expediente) y el petitorio concreto. \\

\textbf{¿Cuál es la diferencia entre un escrito de mero trámite y uno fundado?} & 
El de mero trámite formula un pedido simple sin mayor argumentación; el fundado requiere justificar la causa del pedido (por ejemplo, fuerza mayor) y explicar su relevancia para la causa. \\

\textbf{¿Cuáles son las formas de notificación previstas?} & 
Personal, por cédula (papel o digital) y por nota o ministerio de la ley, esta última los días martes y viernes, salvo las resoluciones del artículo 133 CPCCN que exigen cédula o notificación personal. \\

\textbf{¿Cuándo empieza a correr un plazo procesal?} & 
Nunca el día en que se dicta o notifica la resolución; siempre a partir del primer día hábil siguiente a la notificación efectiva (día de nota: martes o viernes). \\

\textbf{¿Qué es la preclusión procesal?} & 
El cierre definitivo de un acto procesal, ya sea porque se cumplió dentro del plazo o porque el plazo venció sin cumplimiento. Una vez cerrado el acto, no puede reabrirse ni ampliarse (salvo la demanda, antes de notificada). \\

\midrule

\multicolumn{2}{p{\textwidth}}{
\textbf{Síntesis Final:} 

Esta clase recorrió el camino que va desde la norma jurídica abstracta hasta el expediente judicial concreto. Partiendo de una premisa lingüística ---el derecho se trabaja siempre con la palabra--- se llegó a la estructura lógica de toda norma: obligar, permitir y, en cualquiera de los dos casos, sancionar. Sobre esa base se explicó la bifurcación entre el proceso de conocimiento, orientado a establecer los hechos, y el proceso ejecutivo, orientado a verificar la validez formal de un título. Se estudió cómo un conflicto real se transforma en demanda mediante la entrevista con el cliente, la selección estratégica de los hechos probables y la construcción de una verdad procesal que nunca coincide exactamente con la real. El régimen de notificaciones junto con la preclusión cierra la clase con una lección de disciplina procesal: todo plazo mal computado puede costar un derecho.
} \\

\bottomrule
\end{tabularx}
\end{table}

% ==========================================
% FIN DEL DOCUMENTO
% ==========================================

\end{document}
